\section{Introduccion}

% TODO: 3 a 5 parrafos. Guion sugerido:
%   1. Contexto: clasificacion automatica de objetos en linea de produccion.
%   2. Problema: cuanto del desempeno viene del clasificador y cuanto de la
%      representacion de entrada.
%   3. Antecedentes: eigenfaces \parencite{turk1991}, minimos cuadrados como
%      clasificador \parencite{boyd2018}, datasets \parencite{nene1996,muresan2018}.
%   4. Aporte: mismo clasificador lineal sobre tres representaciones, con el
%      motor implementado en NumPy sin bibliotecas de aprendizaje automatico.
%   5. Organizacion del documento.

% Ejemplo de cita para conservar el estilo APA autor-fecha:
% El metodo de eigen-objetos se deriva del trabajo de \textcite{turk1991}.
